% =============================================================
% Sıfırdan İşlemciye — Sayısal Donanım Tasarımı
% Şafak Şahin - Dr. Batuhan Hangün
% Altyapı: Oğuz Ergin'in "Bilgisayar Mimarisi" (prof-oguzergin/bilgisayar-mimarisi)
% kitabının preamble/renk-modu/tikz-sekiller-makrolar/surum-ayarlari
% altyapısından uyarlanmıştır (bkz. _template_ergin/LICENSE).
% =============================================================
\documentclass[12pt, a4paper, twoside, openright]{book}

\input{preamble}
\input{renk-modu}
% =============================================================
% Sıfırdan İşlemciye — Sürüm Ayarları
% Bu dosya sürüm numarası, çalışma modu ve bölüm yönetimini
% merkezi olarak kontrol eder. Yapı, Oğuz Ergin'in
% "Bilgisayar Mimarisi" kitabındaki sistemden uyarlanmıştır.
% =============================================================

% --- Sürüm Bilgisi ---
\newcommand{\kitapsurum}{0.1.0}

% --- Çalışma Modu ---
% Günlük çalışma: sadece etkin bölüm(ler)i derlemek için aşağıdaki satırı aç.
% Tam derleme için yorum satırı yap.
\includeonly{bolum01-sayi-sistemleri/bolum01,bolum02-bool-cebri/bolum02,bolum03-kmap/bolum03}

% --- Bölüm Yönetim Makroları ---
% \bolumekle: Dosya varsa \include eder, yoksa sayacı atlar ve label tanımlar.
%   #1 = dosya yolu (uzantısız), #2 = label adı, #3 = bölüm numarası, #4 = başlık
\makeatletter
\newcommand{\bolumekle}[4]{%
  \IfFileExists{#1.tex}{%
    \include{#1}%
  }{%
    \stepcounter{chapter}%
    \protected@write\@auxout{}{\string\newlabel{#2}{{#3}{}{#4}{chapter.#3}{}}}%
  }%
}
\makeatother


\begin{document}

% --- Ön Sayfalar ---
\frontmatter

% Kapak sayfası şimdilik yok (sonraki bir adımda eklenecek).

\thispagestyle{empty}
\begin{center}
\vspace*{4cm}
{\Huge\bfseries Sıfırdan İşlemciye\par}
\vspace{16pt}
{\Large Sayısal Donanım Tasarımı: Sayı Sistemlerinden Basit Bir CPU Mimarisine\par}
\vspace{24pt}
{\large\bfseries Şafak Şahin - Dr. Batuhan Hangün\par}
\end{center}

\clearpage
\thispagestyle{empty}
\vspace*{\fill}
\noindent
{\small
\textbf{Sıfırdan İşlemciye} --- Sayısal Donanım Tasarımı\\[3pt]
Şafak Şahin - Dr. Batuhan Hangün\\[6pt]
Sürüm \kitapsurum, \the\year\\[6pt]
Kişisel çalışma notu olarak hazırlanmıştır; yayıncı yoktur.\\[6pt]
Bu belgenin LaTeX altyapısı (preamble/renk-modu/tikz-sekiller-makrolar/sürüm-ayarları),
Oğuz Ergin'in \emph{Bilgisayar Mimarisi} kitabından (\url{https://github.com/prof-oguzergin/bilgisayar-mimarisi})
uyarlanmıştır; bkz. o projenin lisansı. Bu belgedeki özgün metin ve teknik içerik yazarına aittir.\\[6pt]
\textit{İlk sürüm, Temmuz 2026}
}
\vspace{2cm}
\cleardoublepage

% İçindekiler
\tableofcontents
\cleardoublepage
\listoffigures
\addcontentsline{toc}{chapter}{Şekil Listesi}
\cleardoublepage
\listoftables
\addcontentsline{toc}{chapter}{Tablo Listesi}
\cleardoublepage
\listofbilginotlari
\cleardoublepage
\listofdetaybilgiler

% --- Ana İçerik ---
\mainmatter

\bolumekle{bolum01-sayi-sistemleri/bolum01}{chap:sayi-sistemleri}{1}{Sayı Sistemleri ve Kodlama}
\bolumekle{bolum02-bool-cebri/bolum02}{chap:bool-cebri}{2}{Bool Cebiri ve Mantık Kapıları}

% --- Arka Sayfalar ---
\backmatter

\chapter*{Terimler Sözlüğü}
\addcontentsline{toc}{chapter}{Terimler Sözlüğü}
% =============================================================
% Terimler Sözlüğü — bölümler ilerledikçe doldurulacak
% =============================================================

\textbf{Taban (Radix)} --- Pozisyonel bir sayı sisteminde kullanılan basamak sayısı; her
basamağın ağırlığı bu sayının kuvvetleriyle belirlenir (ör. ikili için $2$, onaltılık için
$16$).

\textbf{Ardışık Bölme/Çarpma Yöntemi} --- Ondalık bir sayıyı başka bir tabana çevirmek için
kullanılan yöntem; tam sayı kısmı ardışık bölme (kalanlar aşağıdan yukarı okunur), kesirli
kısım ardışık çarpma (taşan basamaklar yukarıdan aşağı okunur) ile çevrilir.

\textbf{Tümleyen (Complement)} --- Bir sayının çıkarma işlemini toplamaya çevirmek için
kullanılan dönüştürülmüş hâli; $(r{-}1)$'in tümleyeni (her basamağı $r{-}1$'den çıkararak) ve
$r$'nin tümleyeni (ona $1$ eklenerek) olmak üzere iki türü vardır.

\textbf{Uçtan-Dolaşan Elde (End-Around Carry)} --- $(r{-}1)$'in tümleyeniyle çıkarma
yapıldığında oluşan taşma bitinin atılmayıp sonucun en düşük anlamlı basamağına eklenmesi
kuralı.

\textbf{İşaret Biti} --- İşaretli bir ikili sayıda en soldaki bit; $0$ pozitif, $1$ negatif
anlamına gelir.

\textbf{Taşma (Overflow)} --- $n$ bitlik iki sayının toplamının $n$ bite sığmaması durumu;
donanımda sabit bit genişliği nedeniyle sonucun bozulmasına yol açar.

\textbf{Normalizasyon} --- Bir kayan noktalı sayıyı $1.xxxx\times2^k$ biçiminde, baştaki
basamak tam olarak $1$ olacak şekilde yazma işlemi; ikilide baştaki $1$ her zaman var olduğu
için saklanmaz (örtük bit).

\textbf{Bias} --- IEEE 754'te üs alanına eklenen sabit kaydırma değeri (tek duyarlıkta
$127$); negatif üsleri de işaretsiz tam sayı olarak karşılaştırılabilir kılar.

\textbf{Bool Değişkeni} --- Yalnızca $0$ veya $1$ değerini alabilen değişken.

\textbf{Mantık Kapısı} --- Bir veya birden çok ikili girişten tek bir ikili çıkış üreten
temel devre elemanı.


\printbibliography[title={Kaynakça}]

\printindex

\end{document}
